\documentclass[11pt]{article}

% ===================== PACKAGES =====================
\usepackage[a4paper,margin=1in]{geometry}
\usepackage{amsmath,amssymb}
\usepackage{enumitem}
\usepackage{fancyhdr}
\usepackage{xcolor} 

% ===================== HEADER & FOOTER =====================
\pagestyle{fancy}
\fancyhf{}
\lhead{CSE 400: Fundamentals of Probability in Computing}
\rhead{Lecture-11 Scribe Submission}
\cfoot{\thepage}

% ===================== CUSTOM COMMANDS =====================
\newcommand{\solution}{
    \vspace{0.5cm}
    \noindent\textbf{\textcolor{blue}{Solution:}}
    \vspace{0.2cm}
    
}

% ===================== TITLE =====================
\title{
    \normalsize School of Engineering and Applied Science (SEAS), Ahmedabad University \\
    \vspace{0.2cm}
    \textbf{CSE 400: Fundamentals of Probability in Computing}\\
    \Large Lecture-19 Scribe Submission
}
\author{} 
\date{}

\begin{document}
\maketitle

\vspace{-2cm}
\begin{center}
    \begin{tabular}{ll}
        \textbf{Name:} & {Jashvi Shah \hspace{2.5in}} \\ [1.5ex]
        \textbf{Enrollment No:} & {AU2440026 \hspace{2.5in}} \\ [1.5ex]
        \textbf{Email:} & {jashvi.s@ahduni.edu.in \hspace{2.1in}} \\ [1.5ex]
    \end{tabular}
\end{center}

\hrule
\vspace{0.5cm}

\section*{CSE400 – Lecture 19 Scribe}

\date{Instructor: Dhaval Patel, PhD \\ Date: March 24, 2026}

\begin{document}

\maketitle

\section*{1. Joint Gaussian Distribution}

\textbf{Definition:}
\[
f_{XY}(x,y)=\frac{1}{2\pi\sigma_X\sigma_Y\sqrt{1-\rho^2}}
\exp\left(-\frac{1}{2(1-\rho^2)}\left[
\left(\frac{x-\mu_X}{\sigma_X}\right)^2+
\left(\frac{y-\mu_Y}{\sigma_Y}\right)^2
-2\rho\left(\frac{x-\mu_X}{\sigma_X}\right)\left(\frac{y-\mu_Y}{\sigma_Y}\right)
\right]\right)
\]

\textbf{Notation:}
\begin{itemize}
\item $\mu_X, \mu_Y$ : means
\item $\sigma_X, \sigma_Y$ : standard deviations
\item $\rho$ : correlation coefficient
\end{itemize}

\section*{2. Example: Joint Gaussian}

\textbf{Given:}
\[
\mu_X=2,\quad \sigma_X^2=1 \Rightarrow \sigma_X=1
\]
\[
\mu_Y=-3,\quad \sigma_Y^2=4 \Rightarrow \sigma_Y=2
\]
\[
\text{cov}(X,Y)=-1
\]

\textbf{Step 1: Correlation}
\[
\rho=\frac{\text{cov}(X,Y)}{\sigma_X\sigma_Y}
=\frac{-1}{1\cdot2}=-\frac{1}{2}
\]

\textbf{Step 2: Substitution}
Substitute into the joint Gaussian formula to obtain the joint PDF.

\section*{3. Marginal PDFs}

\textbf{Definition:}
\[
f_X(x)=\int_{-\infty}^{\infty}f_{XY}(x,y)\,dy,\quad
f_Y(y)=\int_{-\infty}^{\infty}f_{XY}(x,y)\,dx
\]

\textbf{Result:}
\[
X\sim \mathcal{N}(2,1), \quad Y\sim \mathcal{N}(-3,4)
\]

\[
f_X(x)=\frac{1}{\sqrt{2\pi}}e^{-\frac{(x-2)^2}{2}}
\]
\[
f_Y(y)=\frac{1}{2\sqrt{2\pi}}e^{-\frac{(y+3)^2}{8}}
\]

\section*{4. Probability Computation}

\textbf{Step 1:}
\[
P(X<0)=F_X(0)
\]

\textbf{Step 2: Standardization}
\[
Z=\frac{X-2}{1}
\Rightarrow P(X<0)=P(Z<-2)
\]

\textbf{Step 3: Symmetry}
\[
P(Z<-2)=P(Z>2)
\]

\textbf{Final Result:}
\[
P(X<0)=Q(2)
\]

\section*{5. Joint Distribution of N Variables}

\[
\mathbf{X}=
\begin{bmatrix}
X_1 \\ X_2 \\ \vdots \\ X_N
\end{bmatrix}
\]

Joint distribution is described using:
\begin{itemize}
\item Joint PMF
\item Joint PDF
\item Joint CDF
\end{itemize}

\section*{6. Joint Functions}

\textbf{Joint PMF:}
\[
P(X_1=x_1,\dots,X_N=x_N)
\]

\textbf{Joint CDF:}
\[
F(x_1,\dots,x_N)=P(X_1\le x_1,\dots,X_N\le x_N)
\]

\textbf{Joint PDF:}
\[
f(x_1,\dots,x_N)=\frac{\partial^N F}{\partial x_1\cdots\partial x_N}
\]

\section*{7. Independence of Random Variables}

\[
f_{X_1,\dots,X_N}(x_1,\dots,x_N)
=\prod_{i=1}^{N} f_{X_i}(x_i)
\]

\section*{8. Example: Uniform Distribution (3D)}

\textbf{Given:}
\[
f(x_1,x_2,x_3)=
\begin{cases}
C, & |x_i|\le 1 \\
0, & \text{otherwise}
\end{cases}
\]

\textbf{Step 1: Find $C$}
\[
\int_{-1}^{1}\int_{-1}^{1}\int_{-1}^{1}C\,dx_1dx_2dx_3=1
\]
\[
C\cdot 2\cdot 2\cdot 2=1 \Rightarrow C=\frac{1}{8}
\]

\textbf{Step 2: Marginals}
\[
f_{X_1,X_2}(x_1,x_2)=\frac{1}{4}, \quad
f_{X_1}(x_1)=\frac{1}{2}
\]

\textbf{Step 3: Independence}
\[
f_{X_1,X_2,X_3}=\frac{1}{8}
\]
\[
f_{X_1}f_{X_2}f_{X_3}=\frac{1}{2}\cdot\frac{1}{2}\cdot\frac{1}{2}=\frac{1}{8}
\]

\textbf{Conclusion: Independent}

\section*{9. Example: Trivariate Exponential}

\textbf{Given:}
\[
f_{XYZ}(x,y,z)=
\begin{cases}
Ke^{-(ax+by+cz)}, & x,y,z>0 \\
0, & \text{otherwise}
\end{cases}
\]

\textbf{Step 1: Find $K$}
\[
K\int_0^\infty e^{-ax}dx \int_0^\infty e^{-by}dy \int_0^\infty e^{-cz}dz=1
\]
\[
K\cdot \frac{1}{a}\cdot \frac{1}{b}\cdot \frac{1}{c}=1
\Rightarrow K=abc
\]

\textbf{Step 2: Marginals}
\[
f_X(x)=ae^{-ax}
\]

\textbf{Step 3: Independence}
\[
f_{XYZ}=abc e^{-(ax+by+cz)}
\]
\[
f_X f_Y f_Z=abc e^{-(ax+by+cz)}
\]

\textbf{Conclusion: Independent}

\end{document}